\section{The exceptional directions in an individual ramified letter}

% BEGIN P81F1
\paragraph{P81F1. A lattice described by Taylor coefficients.}
Use the actual local data of H80M2, a good coefficient field $L$
with DVR $O$, the invariant parameter $v$ and the degree-two
polynomial $E(v)$ of A81R3. Fix one residue embedding $c$, called
a letter in the source. Put $h=H\ge1$ and $R=O[[v]]$.
Write $D$ for the reduction of the invariant Wach module on the
preceding letter modulo $v$. A81R7 shows that $D$ is free over
$O$ and lifts to an invariant basis. Indeed a $\Delta$-invariant
element in $vN$ equals $v$ times an invariant element, by
cancellation. This identifies $D$ with the original $L_c$ quotient.
In its rational tensor the model acts on the coefficient functions
alone. For $S\subset D$ the native $\mathrm{LamLA}(S)$ is
\[
              \Lambda(S)=E(v)^{-h}(R\otimes_O S).
\]
Here and below only the chosen letter of the product is used.
The coefficient embedding is injective and $E(v)$ is invertible
in the enlarged Laurent coefficient ring used in $\mathrm{LA}$:
$E(v)$ is a unit times $Q$ on this letter (A81R3), and $Q$
reduces to a power of the invertible cyclotomic variable in the
residue Laurent field. These facts justify cancellation there.

At a primitive root $\zeta$ on this letter the source value
$\mathrm{valAt}(y)$ of a lift $y$ of $d\in D$ depends only on
$d$. Indeed the Taylor series of $\varphi y$ is regular, with
value obtained by applying Frobenius to the constant at $1$ on
the preceding letter, by W79H3. Clearing a power of $Q$ multiplies
the first nonzero possible coefficient by the corresponding power
of its leading Taylor coefficient. Dividing back gives exactly
the source value, including when that value is zero. A multiple
of $v$ in the predecessor has zero constant at $1$, giving
well-definedness. Denote this $O$-linear map by $b_\zeta$.
Its image lies in the crystalline period component, hence in
$\mathrm{FilC}_\tau(-h)$ for each associated label $\tau$:
all jumps there are at least $-h$ and all other labels are
unrestricted in $\mathrm{FilC}_\tau$.

For $x\in\Lambda(D)$ write the numerator $E^h x$ in an
integral basis of $D$. Take its Taylor series at each associated
$\zeta$, applying $b_\zeta$ to the basis vectors. Then
$x\in\mathcal N:=\mathrm{NL}$ exactly when every degree-$i$
coefficient lies in $\mathrm{FilC}_\tau(h-i)$ for every such
pair $(\zeta,\tau)$. This is the defining representation of
$\mathrm{NLset}$ using basis lifts. Conversely any allowed finite
presentation has the same Taylor coefficients by the linearity
just proved. The tests are linear, so taking the span in the
definition of $\mathrm{NL}$ preserves them and gives the claimed
description. In particular,
\[
 E^h(R\otimes_O D)\subseteq\mathcal N\subseteq\Lambda(D).
\]
For the first inclusion the numerator is divisible by $E^{2h}$.
Its jets below degree $2h$ vanish since $E$ has a simple zero
at every tested branch. At higher degrees the required filtration
contains the full period component. Multiplication by any element
of $R$ preserves $\mathcal N$, by multiplying Taylor series and
using the decreasing filtration.

Let $\gamma=\gamma_K$ and $u=\chi(\gamma)\in 1+p\mathbb Z_p$
with $v_p(u-1)=1$. This element fixes each $\zeta$ and scales
its logarithmic Taylor variable by $u$. The ratio $\gamma E/E$
is an integral unit. In fact each branch root $a$ of $E$ is
fixed under the substitution on $v$, so $v-a$ divides
$\gamma(v)-a$ integrally; apply the inverse substitution for
equality of the principal ideals. Thus the numerator of $\gamma x$
is obtained by applying this scaling to the old numerator and
multiplying by the unit $(E/\gamma E)^h$. Each Taylor test is
preserved, also for the inverse. Since $\gamma v/v$ is a unit,
$v\mathcal N$ is stable. Finally $v\mathcal N$ is saturated in
$\Lambda(D)$ over $O$. If $\lambda x=v n$ with $n\in\mathcal N$
and $\lambda$ a uniformizer of $O$, evaluating the numerator at
$v=0$ shows $E^h x\in v(R\otimes D)$. Thus $x/v\in\Lambda(D)$.
Its Taylor tests follow by dividing those of $n$ by $\lambda$
in the $L$-linear filtration spaces. This proves saturation on
the actual integral quotient.
% END P81F1

% BEGIN P81F2
\paragraph{P81F2. The rank of the invariant quotient on the two branches.}
The original $\mathrm{IIrest}(S)$ identifies, by P81F1 and the
definition of $\mathrm{InvL}$, with the fixed submodule for
$\gamma$ on
\[
        (\Lambda(S)+v\mathcal N)/(v\mathcal N).
\]
This ambient module is torsion-free over $O$ by saturation.
It is finitely generated, since $v E^{2h}\Lambda(S)\subset
v\mathcal N$ and its successive linear-factor quotients are
finitely generated as follows.

Write $a_1,a_2$ for the two distinct branch values of $v$.
Both lie in the maximal ideal of $O$ and $E=(v-a_1)(v-a_2)$
on this good field. Filter $\Lambda(S)$ by successive
multiplication by the factors of $v E^{2h}$, beginning with
$v$, then $2h$ copies of $v-a_1$, then $2h$ copies of $v-a_2$.
At each step the quotient is generated by a copy of $S$:
evaluation at the current root identifies $R/(v-a)$ with $O$
by division of convergent integral power series (also at $a=0$).
The last term maps to zero in the ambient module above, so the
induced filtration there has graded terms of rank at most
$\operatorname{rank}_O S$.

On the term at zero the scalar on the graded quotient is $1$,
since $E(0)\ne0$ and the substitution fixes zero. At a branch
$a$ let $j$ be the number of previous occurrences of $v-a$.
The multiplier of $\gamma$ on that graded piece is $u^{j-h}$.
Indeed in $X=\zeta\exp(Y)-1$ the series $v(X)-a$ has a simple
zero at $Y=0$. This follows from $v=Z^s$ in A81R3:
$Z'(X)$ is a unit at $X=\zeta-1$, $Z(\zeta-1)\ne0$ and $s$
is prime to $p$. Replacing $Y$ by $uY$ thus multiplies its leading
coefficient by $u$. Every other factor is nonzero at $a$ and
its ratio there is one. The denominator $E^h$ accounts for
the exponent $-h$. Each factor ideal is stable, so this computation
passes to the image filtration, even when the two branch values
coincide modulo the uniformizer.

After taking fraction-field coefficients, a filtration with a
scalar acting on each successive quotient bounds the dimension
of the fixed subspace by the sum of dimensions of the grades
whose scalar is one. This follows by intersecting the filtration
with the fixed space. The cyclotomic unit $u$ has infinite order.
For each branch $0\le j<2h$ only $j=h$ contributes. Including
zero there are precisely three possible grades. The invariant
quotient is a submodule of a finite torsion-free module over the
DVR, hence finite free, and its rank is at most
$3\operatorname{rank}_O S$. This proves the full native
$\mathrm{SecondPropositionFree}$ for each such letter,
$T$ and $S$, with the exact factor $[K:K_0]+1=3$. The same
flag, with its integral divisions, supplies the ring input with
at most three factors of scalar one.
% END P81F2

% BEGIN P81F3
\paragraph{P81F3. The integral factor for the deep jumps.}
At the fixed letter there are two extensions $\tau$ of its
unramified embedding since the good field splits $K$.
The set $\Omega$ in $\mathrm{piOm1}$ is the union of at most
two integer windows. After taking absolute values each consists
of at most $S$ consecutive positive integers at most $h$, since
the endpoint $b_\tau$ is positive and $b_\tau+S\le h$.
The signs distinguish the two deep orientations and do not
affect valuations. For any such interval and any $k\ge1$
the number of multiples of $p^k$ is at most
$\lfloor S/p^k\rfloor+1$, and is zero for $p^k>h$.
Summing in $k$ bounds the sum of the valuations on the interval
by $v_p(S!)+\lfloor\log_p h\rfloor$. Since
$v_p(1-u^i)=1+v_p(|i|)$ for nonzero $i$, the product
$\Pi=\prod_{i\in\Omega}(1-u^i)$ has valuation at most
\[
            \kappa=2\big(S+v_p(S!)+\lfloor\log_p h\rfloor\big)
\]
in the normalization $v_p(p)=1$. A repeated exponent in the two
windows occurs just once in the product, so the same upper
bound applies in that case. For $S=0$ the product is one.
Thus $p^\kappa/\Pi\in O$, giving the exact integral multiplier
required by $\mathrm{KapPi}$ for each pure or mixed letter.

In particular if the first frame proposition supplies a remainder
lattice $S$ of rank at most $2m$, P81F2 bounds its invariant
quotient by $6m$ generators. Any submodule of this finite free
quotient also has at most $6m$ generators over the DVR. This
settles the rank and integral multiplier in the refined count.
The first frame decomposition, lifting its individual coordinates,
and the comparison of the lifted coordinates with Herr classes
have separate content.
% END P81F3

% BEGIN P81F4
\paragraph{P81F4. The two integral changes of basis before Frobenius.}
On the preceding letter let $N=N_{c-1}^\Delta$ over $R=O[[v]]$
with reduction $D=N/vN$. Set $s=(p-1)/2$,
$M=\lceil4h/s\rceil\ge1$, and $\sigma=M-1+v_p((M-1)!)$.
Let $A=R/(v^M)$, $N_M=N/v^MN$ and $c_0=u^s$.
By A81R7 the action of $\gamma$ on $v^jN/v^{j+1}N$ is $c_0^j$.
Over $L$, every $d\in D_L$ has a unique fixed lift to $(N_M)_L$.
Indeed the graded scalars on the positive-degree kernel differ
from one, so $\gamma-1$ is an automorphism there. Correct an
arbitrary lift by its inverse. To control the denominators use
\[
 P(Z)=p^\sigma\prod_{j=1}^{M-1}\frac{Z-c_0^j}{1-c_0^j}\in O[Z].
\]
The product of denominators has normalized valuation
$\sum_{j=1}^{M-1}(1+v_p(j))=\sigma$.
The polynomial kills $v N_M$ by raising through its graded
steps, and $(\gamma-1)N_M\subseteq vN_M$. Hence
$P(\gamma)\widetilde d$ is an integral fixed lift of $p^\sigma d$,
independent of the integral lift $\widetilde d$ of $d$. This
also works for $M=1$. Extending the unique rational fixed lifts
$A_L$-linearly gives an isomorphism
$H:A_L\otimes_L D_L\longrightarrow (N_M)_L$ since its reduction
modulo the nilpotent $v$ is the identity. Its scaled map
$p^\sigma H$ is integral.

For the inverse apply the same calculation to the $A$-linear
dual of $N_M$ with conjugate semilinear action. That dual is
free over $A$ with reduction $D^\vee$, and its positive graded
scalars are the same $c_0^j$ since the action on the reduction
is trivial. Evaluating fixed lifts of a vector and a covector
gives an element of $A_L^\gamma=L$. This equality follows from
the same graded calculation on the positive degrees over $L$
(all $c_0^j-1$ are nonzero there). Their evaluation is therefore
the constant pairing of their reductions. The fixed covectors
are precisely the coordinate functionals of $H^{-1}$, so
$p^\sigma H^{-1}$ also preserves the integral lattices. Both
scaled maps are semilinearly equivariant and natural for all
maps induced by lattice morphisms: an equivariant map sends a
fixed lift to the unique fixed lift of the image of the constant.
% END P81F4

% BEGIN P81F5
\paragraph{P81F5. The full comparison frame for a ramified letter.}
Let $R_c=O[[v]]$ on the output letter and put
$J=(v^M E(v)^{4h})$. Frobenius gives a scalar map
$A\to R_c/J$ because $\varphi(v)=v E^s$ times a unit
by A81R3 and $sM\ge4h$. The invariant part $P_c$ of the
Frobenius span is free on the Frobenius images of the preceding
invariant basis: the images generate the Frobenius span before
invariants and are independent after inverting $Q$ by the two
Wach height inclusions. Hence their invariant coefficients are
exactly in $R_c$. Multiplication by $E^h$ identifies the two
native quotients $\mathrm{QLamLA}$ and $\mathrm{QLamP}$ with,
respectively, the constant $D$ module and $P_c$ modulo $J$.
Indeed their two denominators before multiplication are
$v^M E^{3h}$ times the respective integral modules.

Extend $p^\sigma H$ and $p^\sigma H^{-1}$ from P81F4 along the
Frobenius scalar map above. In the physical module send each
basis vector to its Frobenius translate. This constructs the
required comparison maps $\alpha,\beta$ on the two quotients,
composed with the quotient projections on their native domains.
Their two composites are $p^{2\sigma}$ by P81F4, and both
are linear over the invariant coefficient ring. Frobenius
commutes with the cyclotomic action and with every induced
lattice map. Multiplication by $E^h$ introduces the same unit
$E^h/\gamma(E^h)$ on both numerator modules in the cyclotomic
action. Consequently the maps intertwine the actions and the
maps of lattices in both directions, including the relation
$\mathrm{LAover}$ on finite tensor presentations.

It remains to identify the terms in the primitive Taylor comparison.
Write an arbitrary numerator as a sum of coefficient functions
times classes $d\in D$ and choose lifts of their images under
$p^\sigma H$ in $N$. At the root $1$ on the preceding letter
each lift has constant equal to the period value of $p^\sigma d$
and is fixed by $\gamma$ modulo $v^M$. In the regular period
Taylor series of W79H3 the variable $Y$ is multiplied by $u$;
$v^M$ vanishes to order $sM$ there. Thus all coefficients of
degrees $1$ through $sM-1$ vanish, since $u^j\ne1$ in these
degrees. At a primitive root on the next letter, Frobenius
applies the constant Frobenius and replaces $Y$ by $pY$.
Consequently the image numerator has Taylor coefficients below
$4h$ exactly those of the original functions multiplied by
$p^\sigma b_\zeta(d)$, as in P81F1. Replacing a representative
of the image changes the numerator by $J P_c$, which vanishes
to order at least $4h$ there by regularity. Multiplication by
$Q^k$ shifts agreement through every degree $i<k+4h$, because
$Q$ has a simple zero there. This is precisely the remaining
coefficient identity in $\mathrm{IsFrame}$, for every permitted
clearing exponent $k$ and every presentation. For objects outside
the crystalline bounded-jump test take the zero maps on their
respective domains. All functorial requirements test bounded
crystalline objects on both sides. This gives the full
$\mathrm{FrameA}$ on these fields, simultaneously for the lattices
tested at the letter.
% END P81F5

% BEGIN P81F6
\paragraph{P81F6. From one invariant letter to a finite cohomology class.}
Write $I_j$ for the physical vectors $w$ in $E^{-h}P_c$ with
$(\gamma-1)w\in v^j N^\Delta$. Compare them modulo $v^jN^\Delta$
by intersection with that submodule. The map from the depth-$M$
quotient to the depth-one quotient is injective. Indeed if a
representative is in $vN$, then for $1\le j<M$ its term in
$v^jN^\Delta/v^{j+1}N^\Delta$ is fixed as soon as it lies in
$v^jN^\Delta$. There the scalar is $u^{sj}\ne1$ on a free
$O$-module, so the term vanishes. Induction gives depth $M$.
The image contains $p^\sigma$ times the whole depth-one quotient.
For $w\in I_1$ apply the integral cyclotomic correction of A81R7
to $(\gamma-1)w\in vN^\Delta$. It supplies $k\in vN^\Delta$
with $(\gamma-1)(p^\sigma w-k)\in v^MN^\Delta$.
Project $k$ to letter $c$, which commutes with $\gamma$. The
adjustment stays in $E^{-h}P_c$ because $N_c^\Delta\subseteq
E^{-h}P_c$ by the height inclusions.

For $w\in I_M$ solve $(\varphi-1)z=(\gamma-1)w$ in
$v^M N^\Delta$ on the full set of letters by A81R7. This
gives a Herr pair. If two such first entries differ by
$a\in v^M N^\Delta$, write $a=(\varphi-1)b$ with $b$ in that
same module. The difference of the second entries is exactly
$(\gamma-1)b$, by commutativity and injectivity of $\varphi-1$
on this module. Thus the two classes agree. This gives an
$O$-linear class map on the depth-$M$ quotient. It commutes
with morphisms by the Witt cocycle formula of LWR5, or directly
by applying the morphism to the primitive in that formula.

Its values are finite classes. Indeed after restriction to
$K_1=K_0(\mu_p)$ the cyclotomic kernel and procyclic generator
are the same as before. The Witt formula restricts to the same
pair on these two subgroups. On the actual rings use the old
Wach columns with their original height inclusions after
restriction: the action of this subgroup preserves the ideal $(X)$ and is the
identity on the module modulo $X$. We have $Q^{2h}w\in N$ because $E\sim Q$ and
$Q^hP_c\subset N$, and $z\in v^M N\subset XN$.
The bounded-pair proof B81L2, increasing its height to $2h$,
constructs a crystalline primitive over $K_1$. For the rational
reduction of the integral class, corestriction after restriction
is multiplication by $[K_1:K]$. Transfer preserves finiteness
by its cochain formula and the degree is a unit. Thus the
original class is finite over $K$. All three clauses of the
native $\mathrm{ClassMap}$ and both of $\mathrm{InvLift}$ now
hold on these lattices.
% END P81F6
